\chapter{Topic Modeling}\label{sec-topic-modeling}

\textbf{Learning Objectives}

\emph{Overview starts}

\begin{itemize}
\tightlist
\item
  Understanding what a topic model is.
\item
  Understanding the workings of latent Dirichlet allocation
  (LDA) topic modeling.
\item
  Learning how to apply LDA topic modeling.
\item
  Understanding what a structural topic model (STM) is.
\item
  Extracting information from STMs.
\end{itemize}

\emph{Overview ends}

\section{Identifying latent themes in text
data}\label{identifying-latent-themes-in-text-data}

According to the Oxford Dictionary, a topic is ``a subject that you
talk, write, or learn about''. Or in other terms, it can be understood as
a theme. For instance, a newspaper may cover several topics, such as
politics, crime, lifestyle, the economy, and sport. Each of these topics is
characterized by (combinations of) specific words, something that readers
recognize intuitively. For our tablet dataset and marketing, topics could be the characteristics of certain products,
comparisons of brands, evaluations of and sentiments about products and
services (such as, value for money, product quality, or customer
service), to name but a few.

But what exactly constitutes a model? Depending on the field, it could refer to many elements of an area of research, however, the idea of a model for marketing is linked
with disciplines such as data science and statistics. The
latter deals with data-driven models such as linear regression, which
summarize or describe data sets, ultimately aiming to predict new data
points. Statistical models try to explain an idealized and simplified
reality through mathematics (Micceri 1989). In
statistics, mathematical formulas can be represented by beautiful
graphs and curves that are extremely precise, but real data do not correspond to those shapes. Complex realities are thus transformed
into idealized, simplified mathematical representations. Perfect
straight lines or circles will never describe real data exactly; these
shapes do not exist in the real world. However, the closer the representations
are to a mathematical model, the better we can compare different
data sets. We do this because we need to establish benchmarks for comparison.
Comparison is not just a theoretical exercise, but a practical tool enabling
decision-making, helping marketers to evaluate strategies, forecast
trends, and understand customer behaviors more effectively.

The generic and somewhat pompous definition of topic modeling, is that it is an
unsupervised statistical model used in text analytics which clusters
words into topics without predefined classes, revealing latent themes
within large text corpora. This technique is essential for analyzing
customer reviews or social-media discourses, helping us to detect
trends, sentiments, and emerging issues without reviewing each
single document manually. Typical applications in marketing include, the analysis
of social-media posts or online product reviews (e.g., Tirunillai and
Tellis 2014; Büschken and Allenby 2016), but also purchasing histories
(e.g., Jacobs, Donkers, and Fok 2016), or online browsing
(e.g., Schröder et al. 2019; Trusov, Ma, and Jamal 2016). For an
extensive overview of the application of topic modeling in marketing see
Reisenbichler and Reutterer (2019). From a practical perspective, topic
modeling can inform marketing strategies by identifying the focus of
customers, such as the features of products or service quality. For
instance, in a data set of online hotel reviews, topics such as
`cleanliness' or `service' can be found and further analyzed for
sentiment, frequency, and trends over time, providing insights to help address
the concerns of customer. Determining the optimal number of
topics (K) is a critical step in topic modeling. It is often approached
using educated guesses informed by existing knowledge or systematic
methods like coherence scores. Choosing the right (K), alongside the
interpretation of the words that characterize topics, should, if the practitioners are skilled enough, lead to a blend
of art and science.

Some commonly used topic-modelling techniques used in marketing are the
latent Dirichlet allocation (LDA; Blei, Ng, and Jordan 2003),
structured topic modeling (STM; Roberts, Stewart, and Airoldi 2016),
correlated topic modeling (CTM; Blei and Lafferty 2006),
non-negative matrix factorization (NMF; Lee and Seung 2000), latent
semantic analysis (LSA; Landauer, Foltz, and Laham 1998), or dynamic
topic models (DTM; Blei and Lafferty 2006). For a more complete
list, see Jelodar et al. (2019), but in this chapter we will focus on
topic modeling using LDA and STM.

\section{Latent Dirichlet allocation}\label{latent-dirichlet-allocation}

One of the most studied topic-modeling algorithms is the LDA. The algorithm
aggregates words into topics based on their distribution across
documents, uncovering the `hidden' topics that characterize the corpus.
In practice, we have a set of documents which contain topics
characterized by specific words. The topic exists at a
conceptual level, like a \# (hashtag) which describes the content of a
tweet or an Instagram image. The intuition behind LDA as a
generative model is that writing a document or a review of a product is
driven by themes someone wants to share adds them to the
document. In our data set concerning tablets, one obvious topic might be that of the
characteristics of a device, for example, the screen. Is it bright?
Is it responsive? Or we might consider the connectivity. Is it fast? 5G, Bluetooth
or WiFi? Screen and connectivity are themes or topics. These are
identified by words that are specific to the topic, but also by words which are shared by
multiple topics. In a review, as in any written document, some
themes will be more prevalent than others. If the customer service fell
short of a person's expectations, the theme will emerge. And the words
which characterize it would be more prevalent for this than for other
topics. Let us consider a review which talks about \textbf{connectivity},
(in \(15\%\) of the words contained in the document), about the
\textbf{screen}, (\(25\%\)), \textbf{customer service} (\(50\%\)), and
\textbf{value for money} (\(10\%\)). If we represent this as a
probability distribution over the review, it looks like a vector
{[}0.15, 0.25, 0.50, 0.1{]}. Note that the sum of the topics must be
\(1\).

\begin{verbatim}

Attaching package: 'ggplot2'
\end{verbatim}

\begin{verbatim}
The following object is masked from 'package:NLP':

    annotate
\end{verbatim}

\begin{figure}

\centering{

\includegraphics{7.Chapter_files/figure-pdf/fig-topics_review-1.pdf}

}

\caption{\label{fig-topics_review}Topic Distribution over Review}

\end{figure}%

Each review has a unique distribution of topics, but this is not
all. Each topic is characterized by words, for example, customer
service contains words with the following proportions: bad
(\(0.03\%\)), back (\(0.02\%\)), never (\(0,015\%\)), and so on. Each of
these words in the vocabulary, i.e.,~the whole words used in the tablet-review corpus, is distributed with a relative
frequency over that topic {[}\ldots, 0.02, \ldots, 0.03,
\ldots, 0.015, \ldots{]}. This proportion is the same within
the scope of that topic, no matter which review you take from the
corpus. So, there are two distributions: a mixture of topics from
documents, and a mixture of words from topics. Because these are
all probabilities, we have a probabilistic model here. The single vectors described in previous chapters are characteristic of a
multinomial distribution, but a conjugate prior
distribution over a multinomial distribution is characteristic of the
Dirichlet distribution, hence the name of the algorithm.

To summarize, the LDA model is based on two assumptions: \(K\)
hidden topics are distributed across the entire collection of reviews,
and each review can be represented as a probabilistic distribution of
these topics. Each topic is, in turn, a probabilistic distribution of
words from the data set's vocabulary.

Let us demystify this algorithm further and understand how it functions. To see how LDA is implemented, look at Figure~\ref{fig-lda} (Image
produced by the authors, inspired by Blei, Ng, and Jordan
2003).

\begin{figure}

\centering{

\includegraphics[width=0.9\textwidth,height=\textheight]{figure/chapter7/lda.png}

}

\caption{\label{fig-lda}Latent Dirichlet Allocation}

\end{figure}%

Let us define the terms that appear in the image, then
detail how the algorithm works.

\(M\) is the fixed total number of documents. In a specific corpus
\(D\), documents are indexed by \(i \in \{1,...,M\}\)

\(N\) is the number of words in a given document \(i\). \(N_i\)
represents how many words document \(i\) holds, i.e., it is~the length of the
document.

\(\alpha\) is the (hyper)parameter of the Dirichlet prior to the distribution of topics by document. This describes how topics are
distributed in each document.

\(\beta\) is the (hyper)parameter of the Dirichlet prior to the
word distribution by topic. This describes how words are distributed
inside each topic.

\(\theta_i\) the topic distribution for document \(i\), modeled as a
random variable characterized by a Dirichlet distribution
\(\theta_i \sim \mathcal{D}(\alpha)\).

\(\phi_k\) is the word distribution for topic \(k\), where
\(k \in \{1,...,K\}\), modeled as a random variable characterized by a
Dirichlet distribution \(\phi_k \sim \mathcal{D}(\beta)\)

\(z_{ij}\) is the topic for the \(j^{th}\) word in the \(i^{th}\) document.
The topic is \(z_{i,j} \sim Multinomial(\theta_{i})\).

\(w_{ij}\) is a specific word, \(j\), in document \(i\). The words are
\(w_{i,j} \sim Multinomial(\phi_{z_{i,j}})\)

A multinomial
distribution with one trial is known as a categorial distribution.

The algorithm has a two-step process. The first part is
called the generative model (GM), which describes how documents are
created from a set of latent topics. The second step is called the inference
or estimation procedure (EP) and involves estimating the parameters of
the GM, for example, the topic distributions for each document and the
word distributions for each topic from the documents. It is through this
that we ``learn'' the hidden structure of the topic.

The first step describes how LDA assumes the data are generated and
consist of 5 layers. In layer I in Figure 7.2, a Dirichlet distribution characterized by
parameter \(\alpha\) is used to create the distribution of topics over
each document. If \(\alpha <1\), each document will tend to contain a
few dominant topics (i.e., the distribution is sparse). If
\(\alpha > 1\), the document will tend to contain topics which are distributed
more evenly (Hagg et al. 2022). In layer II, parameter \(\theta\) of a
multinomial distribution, determines the shape of the distribution of
topics over each document. In other words, \(\alpha\) shapes the
multinomial \(\theta\), so each \(\theta_d\) will be different from
document to document, as shown in layer III. In IV we have another
Dirichlet distribution with parameter \(\beta\), this time for all the
words of the vocabulary. If \(\beta<1\), a topic will tend to be
dominated by a smaller set of characteristic words. Starting with the
first word of the first document, to the last word in the last document,
the algorithm will randomly ``assign'' each word to a
topic, based on \(\theta_d\). In layer V, the topic's word distribution
\(\phi_k\) is shared across documents, but be aware that each document's
topic mixture \(\theta_d\) is different. Because each topic \(k\) has a
separate \(\phi_k\), it defines its own ``theme'',characterized by
specific words.

The second step (EP) uses methods such as Gibbs sampling (our focus
here) or variational inference (VI) to estimate \(\theta_d\) for each
document and \(\phi_k\) for each topic from the corpus. We do
not know the true value of these parameters, and we can only observe the
words in each document, but we aim to infer the latent topic assignments
and distributions. In Gibbs sampling, each word's topic assignment is
re-sampled, based on the current assignment of all the other words, until
a stable words-over-topics posterior distribution is reached. Over many
repetitions, the estimates for \(\phi_k\) and \(\theta_d\)
converge to reach a stable state. In this context, ``converge'' means that
the model parameters (\(\phi_k\) and \(\theta_d\)) reach a point where
their distribution stops changing significantly. When Gibbs converges,
we have an estimate for each document's mix of topics \(\theta_d\) and how
the words are distributed within each topic \(\phi_k\).

It is important to use a practical example to understand the underlying
mechanism of Gibbs sampling, once we have looked at that we will provide the mathematical
formulas. Let us use a simplified corpus made up of 4 documents that
resemble the data set in this book:

\(Doc_1\): This tablet has a bright screen.\\
\(Doc_2\): The RAM memory of the tablet is sufficient.\\
\(Doc_3\): The tablet connects to WiFi.\\
\(Doc_4\): The tablet connects to WiFi and has a superb screen.

To further simplify this, we assume to have only two topics, \(T_1\) and
\(T_2\) which represent two themes, ``hardware'' for \(T_1\) and
``connectivity'' for \(T_2\). LDA does not assign labels to topics,
and gives them only generic names such as \(Topic_1\), \(Topic_2\), and so on.
We have to figure out the topic name based on the words assigned to it. Whenever
we apply LDA, we must pre-process the dataset, as we did in the
previous chapters for the bag-of-words approach, such as removing punctuation and stop words, or making all words lowercase. This way, our
vocabulary, in alphabetical order, becomes: \{bright, connects,  memory, screen, sufficient, superb, tablet, wifi\}. Initially, LDA assigns the words
\textbf{randomly} to one of the two topics.

\(Doc_1\): tablet(\(T_1\)), bright (\(T_1\)), screen(\(T_2\))\\
\(Doc_2\): tablet(\(T_2\)), memory(\(T_1\)), sufficient(\(T_2\))\\
\(Doc_3\): tablet(\(T_1\)), connect(\(T_2\)), wifi(\(T_1\))\\
\(Doc_4\): tablet(\(T_2\)), connect(\(T_1\)), wifi(\(T_2\)),
superb(\(T_1\)), screen(\(T_2\))

Based on these, we can create a table to count the
number of occurrences of each word in each topic, see Table~\ref{tbl-tab}. Add reference here

\begin{longtable}[]{@{}lll@{}}
\caption{Term-Topic matrix}\label{tbl-tab}\tabularnewline
\toprule\noalign{}
Term & T1 & T2 \\
\midrule\noalign{}
\endfirsthead
\toprule\noalign{}
Term & T1 & T2 \\
\midrule\noalign{}
\endhead
\bottomrule\noalign{}
\endlastfoot
tablet & 2 & 2 \\
bright & 1 & 0 \\
screen & 0 & 2 \\
memory & 1 & 0 \\
sufficient & 0 & 1 \\
connect & 1 & 1 \\
wifi & 1 & 1 \\
superb & 1 & 0 \\
\end{longtable}

Now we need a document-topic contingency table
(Table~\ref{tbl-documenttopiccounts}), showing the number of words in
each topic assigned to each document.

\begin{longtable}[]{@{}lll@{}}
\caption{Document-Topic
matrix}\label{tbl-documenttopiccounts}\tabularnewline
\toprule\noalign{}
Doc & T1 & T2 \\
\midrule\noalign{}
\endfirsthead
\toprule\noalign{}
Doc & T1 & T2 \\
\midrule\noalign{}
\endhead
\bottomrule\noalign{}
\endlastfoot
1 & 2 & 1 \\
2 & 1 & 2 \\
3 & 2 & 1 \\
4 & 2 & 3 \\
\end{longtable}

We need to count the words in each topic.

\(T_1\): tablet (2) + bright (1) + memory (1) + connect (1) + wifi (1) +
superb (1) = 7 words\\
\(T_2\): tablet (2) + screen (2) + sufficient (1) + connect (1) + wifi
(1) = 7 words

The vocabulary size \(V\) is made up of 8 terms: tablet, bright, screen,
memory, sufficient, connects, wifi, superb.

Now we have arrived at the second step of the algorithm. Following the initial random
assignment, Gibbs sampling is now used to reassign
the words to the existing topics until we obtain convergence, i.e.,~until
the algorithm does not reassign these words to topics anymore. Let us
update the topic assignment for the term ``screen'', which was initially assigned
to Document 1. We must (temporarily) remove one instance of ``screen'' from its current assignment \(T_2\), adjust the counts, and
calculate the conditional probability for each topic. The adjusted count
for \(T_1\) is unchanged, it holds 7 terms while the adjusted count for
\(T_2\) holds 6 terms, you will see a (\(-1\)) in the formula. The result is that we have temporarily removed a term which needs to be freshly assigned to
a topic.

To calculate the conditional probabilities of topic \(k\) given word
\(w\) in document \(d\) we use the formula: Add here a citation

\[
P(T_k|w_d) \propto \left(\frac{n_{k,w}+\beta}{\sum_{w'}(n_{k,w'}+V\beta)} \right) \times \left(\frac{n_{d,k}+\alpha}{\sum_{k'}(n_{d,k'}+K\alpha)}\right) 
\] where

\begin{itemize}
\tightlist
\item
  \(n_{k,w}\) is the count of word \(w\) assigned to topic \(k\),
  excluding the current instance.
\item
  \(n_{d,k}\) is the count of all words in document \(d\) assigned to
  topic \(k\), excluding the current instance.
\item
  \(\sum_{w'}(n_{k,w'}+\beta)\) is the total count of all words assigned
  to topic \(k\) plus \(\beta\) for each word in the vocabulary.
\item
  \(\sum_{k'}(n_{d,k'}+\alpha)\) is the total count of all words in
  document \(d\) plus \(\alpha\) for each topic.
\item
  \(V\) is the vocabulary size.\\
\item
  \(K\) is the number of topics.
\end{itemize}

Be aware that this formula may appear in other books or
articles in a different but similar form, each of them after a similar
outcome.

Assuming hyperparameters \(\alpha = 0.1\) (for document-topic
distributions) and add reference here \(\beta = 0.01\) (for topic-term distributions) we can
calculate the probability of the term ``screen'' to see if it belongs to
\(T_1\) or to \(T_2\).

\(P(T_1 |\text{screen})\) is proportional to:

Term-Topic:
\[\frac{(n_{T_1,\text{screen}} +\beta)}{\text{Total words in} T_1 + V\times \beta} = \frac{0+0.01}{7+8\times 0.01} \]

Document-Topic:
\[\frac{(n_{Doc_1,T_1} +\alpha)}{\text{Total words in }Doc_1 + 2 \times \alpha} = \frac{2+0.1}{3+2\times 0.1} \]

Combined probability for \(T_1\) proportional to:
\(\frac{0.01}{7.08}\times\frac{2.1}{3.2}=0.001608\)

\(P(T_2 |\text{screen})\) is proportional to:

Term-Topic:
\[\frac{(n_{T_2, \text{screen}} +\beta -1)}{\text{Total words in } T_2 + V\times \beta -1} = \frac{2+0.01-1}{7+8\times 0.01-1} \]

Document-Topic:
\[\frac{(n_{Doc_1,T_2}, +\alpha)}{\text{Total words in }Doc_1 + 2 \times \alpha} = \frac{1+0.1}{3+2\times 0.1} \]

Combined probability for \(T_2\) proportional to:
\(\frac{1.01}{6.08}\times\frac{1.1}{3.2}=0.0571\)

After the calculations, we need to compare the results and check where
``screen'' is most likely to belong. We notice that ``screen'' has a
higher chance of being assigned to \(T_2\), because
\(0.0571 > 0.001608\) but the final outcome depends on the many
iterations assigning and reassigning the word to a topic. This is done for each word in each document. In the end, what
changes from the initial random assignment of word to topics, is the
distribution of words depicted in layer IV. The shapes you see in
Figure~\ref{fig-lda} will change. This example has simplified many of the steps,
but it shows the core process of LDA using Gibbs sampling.

\subsection{Applying simple LDA}\label{apply-simple-lda}

To put into practice what has been illustrated, we will use our
tablets data set and apply simple LDA with the default values, so no
fine-tuning will be required, at least at the beginning. For this example, although
there are many R packages which carry out LDA, we will be using
\texttt{seededlda} (Watanabe, Xuan-Hieu, and Watanabe 2022). We will
show later on why the package has this name and why we have chosen it.
Although the computer we used for running the code is capable of
multi-threading, for reproducibility reasons, we have to use a single
thread. This is because the \texttt{set.seed()}, at least while working
with this package, can only ensure reproducibility for random functions
when operating with one thread. If you use multi-threading, the results may vary slightly across runs despite set.seed().
The data set is the one we have already used
in previous chapters of this book. The data pre-processing was performed
according to the package's specifications.

At this point we are trying to find 5 topics, (\(K=5\)) and display the most
frequent terms associated with each topic. Previously, we have seen how
each topic is described by a word distribution, \(\phi_k\), and each
document is characterized by a topic distribution, \(\theta_d\). To help us interpret the results, here we show only the top 10 words in each \(\phi_k\), but
bear in mind that each topic has probabilities for every term in the
vocabulary.

\begin{Shaded}
\begin{Highlighting}[]
\CommentTok{\# Fit an LDA model to the DFM with 5 topics}
\FunctionTok{set.seed}\NormalTok{(}\DecValTok{42}\NormalTok{)}
\NormalTok{lda }\OtherTok{\textless{}{-}} \FunctionTok{textmodel\_lda}\NormalTok{(dfmt, }\AttributeTok{k =} \DecValTok{5}\NormalTok{, }\AttributeTok{verbose =} \ConstantTok{FALSE}\NormalTok{)}
\CommentTok{\# Display terms associated with each topic}
\FunctionTok{terms}\NormalTok{(lda)}
\end{Highlighting}
\end{Shaded}

\begin{verbatim}
      topic1        topic2      topic3     topic4      topic5    
 [1,] "apple"       "browser"   "usb"      "christmas" "support" 
 [2,] "tab"         "button"    "card"     "old"       "service" 
 [3,] "galaxy"      "book"      "keyboard" "happy"     "days"    
 [4,] "samsung"     "content"   "archos"   "perfect"   "customer"
 [5,] "xoom"        "access"    "sd"       "book"      "unit"    
 [6,] "asus"        "cloud"     "windows"  "watch"     "return"  
 [7,] "camera"      "email"     "google"   "gift"      "tried"   
 [8,] "os"          "store"     "flash"    "laptop"    "received"
 [9,] "transformer" "available" "computer" "year"      "problems"
[10,] "thrive"      "page"      "memory"   "wife"      "never"   
\end{verbatim}

The words might give us some indication of the interpretation of the
topics. For example, \texttt{topic\ 4} is related to Christmas gifts,
and \texttt{topic\ 5} to customer service, and so can be labeled by us as
\texttt{gifts} and \texttt{customer\_service}. We can check and see with
which document, represented by a title and a date, a topic is associated. By
examining \(\theta_d\) we know how topics are distributed over each
document. However, the top words give us a hint of how to label and
interpret each topic based on the \(\phi_k\) values.

\begin{Shaded}
\begin{Highlighting}[]
\CommentTok{\# Extract document{-}level variables from the LDA model\textquotesingle{}s data}
\NormalTok{dat }\OtherTok{\textless{}{-}} \FunctionTok{docvars}\NormalTok{(lda}\SpecialCharTok{$}\NormalTok{data)}
\CommentTok{\# Assign the most probable topic to each document based on the LDA model}
\NormalTok{dat}\SpecialCharTok{$}\NormalTok{topic }\OtherTok{\textless{}{-}} \FunctionTok{topics}\NormalTok{(lda)}
\CommentTok{\# Display the first 10 rows of some selected columns}
\FunctionTok{head}\NormalTok{(dat[,}\FunctionTok{c}\NormalTok{(}\StringTok{"Review\_date"}\NormalTok{, }\StringTok{"topic"}\NormalTok{, }\StringTok{"Title"}\NormalTok{)], }\DecValTok{10}\NormalTok{)}
\end{Highlighting}
\end{Shaded}

\begin{verbatim}
   Review_date  topic                                                Title
1     1/8/2012 topic1                                 Love My Kindle Fire!
2    4/10/2012 topic4                    Not the best Tablet or a good buy
3    3/26/2012 topic2            Had it 3 months, would not purchase again
4    6/14/2012 topic4 College Grad - I use it waaaay more than my lap top!
5    5/21/2011 topic4                              Better than I expected.
6    1/29/2010 topic5                            Pick a different computer
7   11/26/2011 topic3                     Great for travel and general use
8     5/6/2012 topic1                         Highly Underestimated Tablet
9    6/29/2011 topic3                          Lenovo ThinkPad X220 Tablet
10   6/21/2011 topic2                              A little disappointment
\end{verbatim}

In the output above, a row corresponds to a document (i.e., a review), the
\texttt{Topic} column shows the most probable topic assigned to that
document. The other two columns, \texttt{Review\_date} and
\texttt{Title}, help us identify documents and we can check how the model
has clustered them by topic. We know that each document has the topic
distribution \(\theta_d\), but for clarity we display only the
topic with the highest probability.

It is possible to increase and decrease the number of topics to understand
better which themes can be found in the text. To speed up the execution
time one can parallelize \texttt{seededlda} by adding a
\texttt{batch\_size} option and enable convergence
\texttt{auto\_iter=TRUE}, for example:
\texttt{textmodel\_lda(dfmt,\ k\ =\ 5,\ batch\_size\ =\ 0.01,\ auto\_iter\ =\ TRUE)}.

\subsection{Seeded LDA.}\label{seeded-lda.}

Let us assume that you, as an analyst, already know the topics (e.g., from
a previous study on a different, but related corpus of text), but you want
LDA to group the words in the topics with increased precision. Seeded
LDA allows for semi-supervised topic modeling, acknowledging that you
might have prior knowledge of your dataset and you can provide a
dictionary that links words to topics which have already been labeled.
These ``seeded'' words provide a starting point for the algorithm,
meaning that the random assignment, as in the first step of the LDA above, becomes ``informed'' by the dictionary. In
order to perform this we need a dictionary function and an \texttt{yml}
file to store the information. Each line defines a topic's label, such as
``hardware'' or ``functionality'', and associates it with a list of
``seed words'' which are strongly related to that topic. The structure
of the file, \texttt{dictionary.yml}, looks like this.

hardware: {[}ram, gb, storage, slot, external, camera, power, wifi,
battery, button{]}

functionality: {[}video*, games, *books, email, social, media, reader,
browser{]}

customer\_service: {[}service, support, customer, return{]}

names: {[}amazon, apple, ipad, xoom, samsung, galaxy, google{]}

experience: {[}screen*, display, resolution, glass*, touch, thin,
light{]}

where the * represents a placeholder that captures plurals or other
inflectional forms. For example, video* could match video or videos, and
*books can capture terms like books, songbooks, or notebooks.

When we employ Seeded LDA, the algorithm uses the seeded words to
initialize (guide) the assignments so that the model places them
(and similar words) in the right topic. Compared to the LDA described above, the random-assignment step is replaced with an
``informed'' assignment based on domain knowledge.

\begin{Shaded}
\begin{Highlighting}[]
\CommentTok{\# Create the dict variable containing the seeded words}
\NormalTok{dict }\OtherTok{\textless{}{-}} \FunctionTok{dictionary}\NormalTok{(}\AttributeTok{file =} \StringTok{"data/dictionary.yml"}\NormalTok{)}

\CommentTok{\# Fit an seeded LDA model to the DFM with 5 topics}
\FunctionTok{set.seed}\NormalTok{(}\DecValTok{42}\NormalTok{)}
\CommentTok{\# Add auto\_iter = TRUE for faster convergence}
\NormalTok{lda\_seed }\OtherTok{\textless{}{-}} \FunctionTok{textmodel\_seededlda}\NormalTok{(dfmt, dict, }\AttributeTok{auto\_iter =} \ConstantTok{TRUE}\NormalTok{)}
\CommentTok{\# Display terms associated with each topic}
\FunctionTok{terms}\NormalTok{(lda\_seed)}
\end{Highlighting}
\end{Shaded}

\begin{verbatim}
      hardware   functionality customer_service names         experience  
 [1,] "camera"   "email"       "support"        "apple"       "display"   
 [2,] "power"    "browser"     "service"        "google"      "light"     
 [3,] "button"   "videos"      "customer"       "samsung"     "christmas" 
 [4,] "usb"      "reader"      "return"         "galaxy"      "resolution"
 [5,] "storage"  "media"       "days"           "xoom"        "happy"     
 [6,] "keyboard" "access"      "unit"           "tab"         "old"       
 [7,] "card"     "content"     "tried"          "os"          "laptop"    
 [8,] "sd"       "book"        "update"         "asus"        "perfect"   
 [9,] "slot"     "store"       "received"       "flash"       "book"      
[10,] "windows"  "cloud"       "problems"       "transformer" "watch"     
\end{verbatim}

Notice that all the topic labels are described nicely
by the seeded words, grouping other pertinent words to those
topics. We might notice that some words in the dictionary file do not
appear as the top words in the results. This might be because those words, although they belong to that topic, might not be so
prevalent in the corpus. You can list more than 10 terms using the
function \texttt{terms(lda\_seed,\ 20)}, for example. Let us assume that
you are aware of the 5 topics above are represented in your data set
(e.g., based on previous studies or observations of other review data)
but you have a hunch, that there may be more topics that might be useful. In this case, you can use seeded LDA with residual
topics, which is a combination of both the seeded and unseeded options. We
demonstrate this below.

\begin{Shaded}
\begin{Highlighting}[]
\CommentTok{\# Fit an seeded LDA model to the DFM with 7 topics}
\FunctionTok{set.seed}\NormalTok{(}\DecValTok{42}\NormalTok{)}
\CommentTok{\# Add auto\_iter = TRUE for faster convergence}
\NormalTok{lda\_res }\OtherTok{\textless{}{-}} \FunctionTok{textmodel\_seededlda}\NormalTok{(dfmt, dict, }\AttributeTok{residual =} \DecValTok{2}\NormalTok{, }\AttributeTok{auto\_iter =} \ConstantTok{TRUE}\NormalTok{)}
\CommentTok{\# Display terms associated with each topic}
\FunctionTok{terms}\NormalTok{(lda\_res)}
\end{Highlighting}
\end{Shaded}

\begin{verbatim}
      hardware   functionality customer_service names         experience  
 [1,] "camera"   "email"       "support"        "apple"       "display"   
 [2,] "power"    "browser"     "service"        "google"      "light"     
 [3,] "button"   "videos"      "customer"       "samsung"     "resolution"
 [4,] "usb"      "reader"      "return"         "galaxy"      "keyboard"  
 [5,] "storage"  "media"       "days"           "xoom"        "bit"       
 [6,] "card"     "access"      "unit"           "tab"         "hold"      
 [7,] "sd"       "content"     "tried"          "os"          "volume"    
 [8,] "slot"     "store"       "update"         "asus"        "pretty"    
 [9,] "keyboard" "book"        "problems"       "transformer" "sound"     
[10,] "windows"  "cloud"       "received"       "honeycomb"   "screens"   
      other1      other2    
 [1,] "christmas" "people"  
 [2,] "book"      "reviews" 
 [3,] "happy"     "hp"      
 [4,] "old"       "touchpad"
 [5,] "gift"      "review"  
 [6,] "perfect"   "anything"
 [7,] "wife"      "never"   
 [8,] "year"      "bad"     
 [9,] "watch"     "worth"   
[10,] "prime"     "buying"  
\end{verbatim}

Note that the two extra topics, called \texttt{other1} and
\texttt{other2}, represent new information. The first one seems to refer to tablets as gifts for Christmas, and the second seems
to be a comparison between tablets and laptops or computers. From a marketing point of
view, the appearance of \texttt{other1} might be interesting for
understanding when (and where) such a product can be advertised, and the
other topic reveals competition with other categories of product.

\subsection{Applying sequential sentence-based LDA with seeded
words}\label{apply-sequential-lda-sentence-based-with-seeded-words}

Sentence-based LDA (see, e.g., Büschken and Allenby 2016) is an
interesting way to tokenize a corpus into documents as sentences. In this,
each sentence in a review represents a document. Several studies have pointed
out that when people leave a review, topics are expressed in
sentences, and rarely with two or more topics in a single sentence.
Achieving this is solved using the
\texttt{corpus\_reshape()} function. In this case we will consider the
corpus as a merged list of sentences. Let us apply sequential LDA and
check the results. The result of the document-feature matrix (DFM),
which is the equivalent of the document-term matrix, will show as many
documents as the number of sentences (i.e., documents), and the sparsity of
such a matrix is very high. Use \texttt{print(dfmt)} to print it out. The
sequential LDA algorithm incorporates the topics discovered for the
previous sentence when modeling the next one, effectively acknowledging
that neighboring sentences in a review often share related subjects. The previous sentence does not force a topic into the
following one, but influences the model's estimation of whether it belongs there. In this case, each
sentence's topic distribution is ``guided'' but not determined by the preceding one, showing a bias towards continuity.

\begin{Shaded}
\begin{Highlighting}[]
\FunctionTok{library}\NormalTok{(seededlda)}
\FunctionTok{library}\NormalTok{(quanteda)}
\FunctionTok{library}\NormalTok{(ggmatplot)}


\NormalTok{corp }\OtherTok{\textless{}{-}} \FunctionTok{corpus}\NormalTok{(prodreviews, }\AttributeTok{text\_field =}\StringTok{"Content"}\NormalTok{) }\SpecialCharTok{|\textgreater{}}
    \FunctionTok{corpus\_reshape}\NormalTok{(}\AttributeTok{to =} \StringTok{"sentences"}\NormalTok{)}

\NormalTok{toks }\OtherTok{\textless{}{-}} \FunctionTok{tokens}\NormalTok{(corp, }\AttributeTok{remove\_punct =} \ConstantTok{TRUE}\NormalTok{, }\AttributeTok{remove\_symbols =} \ConstantTok{TRUE}\NormalTok{, }
               \AttributeTok{remove\_numbers =} \ConstantTok{TRUE}\NormalTok{, }\AttributeTok{remove\_url =} \ConstantTok{TRUE}\NormalTok{)}
\NormalTok{dfmt }\OtherTok{\textless{}{-}} \FunctionTok{dfm}\NormalTok{(toks) }\SpecialCharTok{|\textgreater{}} 
    \FunctionTok{dfm\_remove}\NormalTok{(}\FunctionTok{stopwords}\NormalTok{(}\StringTok{"en"}\NormalTok{)) }\SpecialCharTok{|\textgreater{}}
    \FunctionTok{dfm\_remove}\NormalTok{(}\StringTok{"*@*"}\NormalTok{) }\SpecialCharTok{|\textgreater{}}
    \FunctionTok{dfm\_trim}\NormalTok{(}\AttributeTok{max\_docfreq =} \FloatTok{0.1}\NormalTok{, }\AttributeTok{docfreq\_type =} \StringTok{"prop"}\NormalTok{)}

\CommentTok{\# Enable the sequential algorithm, set gamma to 0.5}
\CommentTok{\# A lower value of gamma would impact the classification of the current sentence.}
\NormalTok{lda\_seq }\OtherTok{\textless{}{-}} \FunctionTok{textmodel\_lda}\NormalTok{(dfmt, }\AttributeTok{k =} \DecValTok{5}\NormalTok{, }\AttributeTok{gamma =} \FloatTok{0.5}\NormalTok{, }
                         \AttributeTok{batch\_size =} \FloatTok{0.01}\NormalTok{, }\AttributeTok{auto\_iter =} \ConstantTok{TRUE}\NormalTok{,}
                         \AttributeTok{verbose =} \ConstantTok{FALSE}\NormalTok{)}
\FunctionTok{print}\NormalTok{(lda\_seq)}
\end{Highlighting}
\end{Shaded}

\begin{verbatim}

Call:
textmodel_lda(x = dfmt, k = 5, auto_iter = TRUE, gamma = 0.5, 
    batch_size = 0.01, verbose = FALSE)

5 topics; 442,621 documents; 60,204 features.
\end{verbatim}

After running the \texttt{textmodel\_lda()} function, we obtain the topics which are displayed in
Table~\ref{tbl-sentence_based_lda2}. Each column represents the words
that the model thinks that are most characteristic of that
topic. By examining these terms, we can analyze, sentence-by-sentence, which concept each topic
encapsulates.

\begin{longtable}[]{@{}lllll@{}}
\caption{Top Terms for Each Topic (Sentence-Based
LDA)}\label{tbl-sentence_based_lda2}\tabularnewline
\toprule\noalign{}
topic1 & topic2 & topic3 & topic4 & topic5 \\
\midrule\noalign{}
\endfirsthead
\toprule\noalign{}
topic1 & topic2 & topic3 & topic4 & topic5 \\
\midrule\noalign{}
\endhead
\bottomrule\noalign{}
\endlastfoot
screen & kindle & tablet & tablet & apps \\
just & fire & one & screen & android \\
device & books & ipad & can & app \\
get & love & product & use & tablet \\
button & amazon & get & ipad & can \\
touch & can & just & like & use \\
one & read & amazon & good & device \\
can & use & first & battery & like \\
like & great & now & great & market \\
user & reading & android & usb & great \\
\end{longtable}

\subsection{Topic visualization}\label{topics-visualization}

Topics, and how they are related to each other, can be visualized with the use of the
LDAVis package (Sievert and Shirley 2014). This R package is an
interactive web-based tool designed for visualizing the output of LDA topic models. The plots show the relationships between
topics, emphasize the importance of each topic, and the top terms
associated with each topic. The tool is mostly self-explanatory and
helps us to interpret the results.

\begin{Shaded}
\begin{Highlighting}[]
\FunctionTok{library}\NormalTok{(LDAvis)}
\FunctionTok{library}\NormalTok{(quanteda)}
\FunctionTok{library}\NormalTok{(stringi)}

\CommentTok{\# Find required quantities}
\NormalTok{json }\OtherTok{\textless{}{-}} \FunctionTok{createJSON}\NormalTok{(}\AttributeTok{phi =}\NormalTok{ lda\_seq}\SpecialCharTok{$}\NormalTok{phi,}
                    \AttributeTok{theta =}\NormalTok{ lda\_seq}\SpecialCharTok{$}\NormalTok{theta,}
                    \AttributeTok{doc.length =} \FunctionTok{rowSums}\NormalTok{(dfmt),}
                    \AttributeTok{vocab =} \FunctionTok{colnames}\NormalTok{(dfmt),}
                    \AttributeTok{term.frequency =} \FunctionTok{colSums}\NormalTok{(dfmt)}
\NormalTok{)}

\FunctionTok{serVis}\NormalTok{(json)}
\end{Highlighting}
\end{Shaded}

\begin{figure}

\centering{

\includegraphics[width=0.9\textwidth,height=\textheight]{figure/chapter7/ldavis2.png}

}

\caption{\label{fig-ldavis2}LDAvis}

\end{figure}%

Figure~\ref{fig-ldavis2} is a static visualization of the topics
generated using LDA. On the left-hand side, we can see a map displaying
the distances between topics. These can be interpreted as their degree of
similarity. When topics are close to each other or
overlapping, they share more terms. The map is generated using
multidimensional scaling, a technique which projects
high-dimensional data into a two-dimensional plane while preserving distances between pairs as much as possible. Each of the 5 topics is
represented by a circle (the marginal topic distribution) whose size is
proportional to the prevalence of the topic in the data set. We can see
that the sizes are similar. The navigation between topics to get more
detail can be done through the controls such as ``Selected Topic'', while
``Top-30 Most Salient Terms'' ranks the words based on a combination
of their frequency and relevance to the topic. The latter can be adjusted
through the parameter \(\lambda\) which controls the balance between
the relevance and exclusivity of terms. When \(\lambda=0\), the focus is on
the terms which are unique to the topic, showing their exclusivity, while
when \(\lambda=1\), the ranking unveils the most frequent terms within the
topic, even if they are common across the whole corpus. \(\lambda\)
allows to us to see the space between the specificity and generality of the terms. It
gives the ratio between the probability of a term occurring within a topic and its marginal
probability across the corpus. The visual representation allows see which terms are combined with which topics quickly.

\section{Structural topic models}\label{structural-topic-models}

A structural topic model (STM; Roberts et al. 2013, 2014) was created with existing topic models in mind (i.e., LDA; Blei, Ng, and Jordan 2003) as well as correlated topic modelling (CTM; Blei and Lafferty 2007). While LDA assumes that topics are
independent, CTM relaxes this assumption by allowing topics to be
correlated. This approach can be suitable for the cases where certain
topics have the tendency to occur at the same time. STM complements to this by allowing users to add extra information, known as metadata addition, about each document to be
analyzed within the topic modeling framework. This procedure helps marketers discover topics and link
them to covariates represented by the related document metadata, because
topical prevalence and topical content can be a function of these
(Roberts, Stewart, and Tingley 2019). Technically, metadata which
influence the popularity of topics are referred to as topical prevalence
covariates, and variables (words) which define the themes of topics are
called topical content covariates.

In this section we will use the R \texttt{stm} package (Roberts,
Stewart, and Tingley 2019) and its functions for textual data and their associated metadata. In our case,
this can be the review, as the text, the associated brand, the review's
date, and it's rating are metadata. We can use the \texttt{quanteda}
package to import the text and metadata, and then pre-process and transform the
document into a document-term matrix. From the \texttt{prodreviews}
variable, which contains the data set, we can select the Content, Brand,
Review\_date, and Rating columns. The function \texttt{prepDocuments()}
uses the package \texttt{tm} to pre-process the data.

The workflow and some of the most important features of the \texttt{stm} package are illustrated in Figure~\ref{fig-stmheuristic}.

\begin{figure}

\centering{

\includegraphics[width=0.9\textwidth,height=\textheight]{figure/chapter7/stmheuristic.png}

}

\caption{\label{fig-stmheuristic}Analysis workflow and stm package
features, Figure from\ldots{} (Roberts, Stewart, and Tingley 2019)}

\end{figure}%

Notice that during the data-ingestion stage, the data set is loaded into R's
memory. Then, the relevant columns are selected, and the text, together
with the metadata, is pre-processed. An important part is the selection
of the brands we want to focus on. We will use Amazon and Apple and aim
to extract background information regarding them.

\begin{Shaded}
\begin{Highlighting}[]
\FunctionTok{library}\NormalTok{(stm)}
\FunctionTok{library}\NormalTok{(tm)}
\FunctionTok{library}\NormalTok{(dplyr)}
\FunctionTok{library}\NormalTok{(lubridate)}

\CommentTok{\# Preprocessing}
\FunctionTok{options}\NormalTok{(}\AttributeTok{stringsAsFactors =} \ConstantTok{FALSE}\NormalTok{)}

\CommentTok{\# Read product reviews from a CSV file into a dataframe}
\NormalTok{prodreviews }\OtherTok{\textless{}{-}} \FunctionTok{read.csv}\NormalTok{(}\StringTok{"data/Product Reviews.csv"}\NormalTok{)}

\CommentTok{\# Select the desired columns}
\NormalTok{data }\OtherTok{\textless{}{-}}\NormalTok{ prodreviews[, }\FunctionTok{c}\NormalTok{(}\StringTok{"Content"}\NormalTok{, }\StringTok{"Rating"}\NormalTok{, }\StringTok{"Brand"}\NormalTok{, }\StringTok{"Review\_date"}\NormalTok{)]}

\CommentTok{\# Transform the character content of Review\_date into date type}
\NormalTok{data}\SpecialCharTok{$}\NormalTok{Review\_date }\OtherTok{\textless{}{-}} \FunctionTok{mdy}\NormalTok{(data}\SpecialCharTok{$}\NormalTok{Review\_date)}

\CommentTok{\# Select the year 2012}
\NormalTok{data }\OtherTok{\textless{}{-}}\NormalTok{ data }\SpecialCharTok{\%\textgreater{}\%} \FunctionTok{filter}\NormalTok{(}\FunctionTok{year}\NormalTok{(Review\_date) }\SpecialCharTok{==} \DecValTok{2012}\NormalTok{)}

\CommentTok{\# Use the strftime to transform the date into a number (day of the year)}
\NormalTok{data}\SpecialCharTok{$}\NormalTok{Review\_date }\OtherTok{\textless{}{-}} \FunctionTok{as.integer}\NormalTok{(}\FunctionTok{strftime}\NormalTok{(data}\SpecialCharTok{$}\NormalTok{Review\_date, }\AttributeTok{format =}\StringTok{"\%j"}\NormalTok{))}

\CommentTok{\# Select the relevant rows (Amazon and Apple)}
\NormalTok{data }\OtherTok{\textless{}{-}} \FunctionTok{subset}\NormalTok{(data, Brand }\SpecialCharTok{==} \StringTok{"Amazon"} \SpecialCharTok{|}\NormalTok{ Brand }\SpecialCharTok{==} \StringTok{"Apple"}\NormalTok{)}

\CommentTok{\# Set seed for reproducibility}
\FunctionTok{set.seed}\NormalTok{(}\DecValTok{42}\NormalTok{)}

\CommentTok{\# Process the text data}
\NormalTok{processed }\OtherTok{\textless{}{-}} \FunctionTok{textProcessor}\NormalTok{(}\AttributeTok{documents =}\NormalTok{ data}\SpecialCharTok{$}\NormalTok{Content, }\AttributeTok{metadata =}\NormalTok{ data)}
\end{Highlighting}
\end{Shaded}

After this, we need to prepare the corpus by linking the text
with metadata using the \texttt{prepDocuments()} utility function. In
building the model, we can set a threshold (e.g., through the option
\texttt{lower.thresh}) for the minimum number of documents in which a
word must appear. Note this has an influence on the rare words. The
threshold helps to reduce noise by excluding terms that are occur infrequently
and might not be helpful for the topic modeling. After we have chosen a
threshold, we can also extract some important variables, such as the
documents, as a large list in \texttt{docs}, the vocabulary as a
character vector in \texttt{vocab}, and the metadata, as a table, in
\texttt{meta}.

\begin{Shaded}
\begin{Highlighting}[]
\CommentTok{\# Associate text with metadata}
\NormalTok{out }\OtherTok{\textless{}{-}} \FunctionTok{prepDocuments}\NormalTok{(}\AttributeTok{documents =}\NormalTok{ processed}\SpecialCharTok{$}\NormalTok{documents, }
                     \AttributeTok{vocab =}\NormalTok{ processed}\SpecialCharTok{$}\NormalTok{vocab,}
                     \AttributeTok{meta =}\NormalTok{ processed}\SpecialCharTok{$}\NormalTok{meta, }\AttributeTok{lower.thresh =} \DecValTok{15}\NormalTok{)}
\end{Highlighting}
\end{Shaded}

\begin{verbatim}
Removing 15119 of 17310 terms (32987 of 483031 tokens) due to frequency 
Removing 3 Documents with No Words 
Your corpus now has 10235 documents, 2191 terms and 450044 tokens.
\end{verbatim}

\begin{Shaded}
\begin{Highlighting}[]
\CommentTok{\# Extract variables}
\NormalTok{docs }\OtherTok{\textless{}{-}}\NormalTok{ out}\SpecialCharTok{$}\NormalTok{documents}
\NormalTok{vocab }\OtherTok{\textless{}{-}}\NormalTok{ out}\SpecialCharTok{$}\NormalTok{vocab}
\NormalTok{meta }\OtherTok{\textless{}{-}}\NormalTok{ out}\SpecialCharTok{$}\NormalTok{meta}
\end{Highlighting}
\end{Shaded}

Next we can estimate a topic STM model by passing the output to the
\texttt{stm()} function. \texttt{K} represents the number of topics we
want to model. The \texttt{prevalence} option, refers to how the proportions of topics
vary across the documents. The covariates can be encoded through
a formula which accommodates multiple covariates and variables that can
be factorial or continuous. In this example, prevalence is a function
of the \texttt{Brand} variable, which generally includes many brands,
but in our case only two. In the model you can set the maximum number of
iterations through the \texttt{max.em.its} option. Spectral
initialization is a great optimization tool for getting the same
results each time we apply the \texttt{stm()} function (Roberts,
Stewart, and Airoldi 2016). Be aware, that this might terminate the
model before convergence is reached. Check the package's manual for more
options (Roberts, Stewart, and Tingley 2019).

\begin{Shaded}
\begin{Highlighting}[]
\CommentTok{\# Fitting the example model}
\NormalTok{productsFit }\OtherTok{\textless{}{-}} \FunctionTok{stm}\NormalTok{(}\AttributeTok{documents =}\NormalTok{ out}\SpecialCharTok{$}\NormalTok{documents, }\AttributeTok{vocab =}\NormalTok{ out}\SpecialCharTok{$}\NormalTok{vocab,}
                       \AttributeTok{K =} \DecValTok{5}\NormalTok{, }\AttributeTok{prevalence =}\SpecialCharTok{\textasciitilde{}}\NormalTok{ Brand, }
                       \AttributeTok{max.em.its =} \DecValTok{75}\NormalTok{, }\AttributeTok{data =}\NormalTok{ out}\SpecialCharTok{$}\NormalTok{meta,}
                       \AttributeTok{init.type =} \StringTok{"Spectral"}\NormalTok{, }\AttributeTok{verbose =} \ConstantTok{FALSE}\NormalTok{)}
\end{Highlighting}
\end{Shaded}

Now we can try to understand and interpret the model estimates, and use the functions. \texttt{labelTopics()} will display the
most likely words associated with the topics. By default, it
uses multiple methods to identify representative words for each topic
(highest prob, FREX, lift, and ccore). In the example below, the output
from \texttt{labelTopics()} shows the results for all the four methods you
can see under ''Topic's Top Words'', and allows you to compare the
different ways topics are represented. The words for each topic
vary depending on the method used. To clarify the effects of using this
methods, \textbf{highest prob} lists the most frequent words within each
topic and also some common words, reducing specificity.
\textbf{FREX}(\textbf{FR}equency and \textbf{EX}clusivity) identifies
words which are both frequent and exclusive to a topic, making them more
distinctive. The words that rank highly in FREX occur frequently in the topic
and are unique to it, in a similar way to TF-IDF \textbf{Lift}
identifies words which are common within and specific to one topic, compared to those which can be found across
the entire corpus. It can be helpful in highlighting unique terms for smaller topics.
\textbf{Score} uses a combination of the log frequency inside one topic and the
penalizing of the word's log frequency across all topics; the goal is to
find words which are common in a specific topic but uncommon elsewhere.

\begin{Shaded}
\begin{Highlighting}[]
\CommentTok{\# Find which words belong to which topic}
\NormalTok{(lT }\OtherTok{\textless{}{-}} \FunctionTok{labelTopics}\NormalTok{(productsFit))}
\end{Highlighting}
\end{Shaded}

\begin{verbatim}
Topic 1 Top Words:
     Highest Prob: amazon, kindl, problem, work, one, get, fire 
     FREX: return, custom, credit, servic, call, password, told 
     Lift: non, paperweight, spoke, worthless, refund, rep, south 
     Score: non, custom, return, problem, credit, call, amazon 
Topic 2 Top Words:
     Highest Prob: app, devic, amazon, fire, use, screen, tablet 
     FREX: browser, volum, cloud, button, content, googl, speaker 
     Lift: clever, launcher, onboard, robust, wallpap, root, upsid 
     Score: clever, android, app, devic, browser, button, volum 
Topic 3 Top Words:
     Highest Prob: love, great, use, can, one, game, get 
     FREX: game, kid, play, daughter, wife, love, gift 
     Lift: toss, hubbi, girl, bus, appoint, grandson, granddaught 
     Score: love, toss, game, great, christma, play, gift 
Topic 4 Top Words:
     Highest Prob: ipad, tablet, appl, use, new, app, like 
     FREX: appl, ipad, iphon, con, camera, itun, compar 
     Lift: retina, con, lte, sim, facetim, dictat, imac 
     Score: con, ipad, appl, retina, iphon, camera, tablet 
Topic 5 Top Words:
     Highest Prob: kindl, fire, book, read, love, use, purchas 
     FREX: color, kindl, book, fire, librari, read, enjoy 
     Lift: proud, sun, sunglass, anti-glar, paperback, lamp, beach 
     Score: kindl, fire, book, proud, read, love, color 
\end{verbatim}

Let us better understand the model's results by using visualizations. We
can explore the words associated with each topic by plotting them. The
top three words associated with the list from the highest prob method 
appear in Figure~\ref{fig-stm4}.

\begin{Shaded}
\begin{Highlighting}[]
\FunctionTok{plot}\NormalTok{(productsFit, }\AttributeTok{type =} \StringTok{"summary"}\NormalTok{, }\AttributeTok{xlim =} \FunctionTok{c}\NormalTok{(}\DecValTok{0}\NormalTok{, }\FloatTok{0.8}\NormalTok{))}
\end{Highlighting}
\end{Shaded}

\begin{figure}[H]

\centering{

\includegraphics{7.Chapter_files/figure-pdf/fig-stm4-1.pdf}

}

\caption{\label{fig-stm4}}

\end{figure}%

According to the plot displayed in Figure~\ref{fig-stm4}, topic 3
(characterized by words like ''love'', ''great'', and ''use'') represents a major
topic, followed by 5, 1, and 2. Topic 4 is a relatively minor topic.

In case we want to see which documents are associated with which
topics we can use the \texttt{findThoughts()} function. This retrieves the
most representative documents for a topic, based on their
proportion in the STM model. By examiningthese documents, we have
the chance to better understand the content of a topic. We have
shortened the documents to 200 characters as a dimensionality reduction,
and checked the documents associated with topics 4 and 5. The documents
selected for each topic are the ones which are most strongly aligned with it. The example code below uses  the \texttt{labelTopics()}
function, which helps us explore the documents strongly associated
with the words identified by the various labeling methods.

\begin{Shaded}
\begin{Highlighting}[]
\CommentTok{\# Select the first 200 text characters, to simplify}
\NormalTok{shortdoc }\OtherTok{\textless{}{-}} \FunctionTok{substr}\NormalTok{(meta}\SpecialCharTok{$}\NormalTok{Content, }\DecValTok{1}\NormalTok{, }\DecValTok{200}\NormalTok{)}

\CommentTok{\# Find 2 documents associated with topic 4}
\NormalTok{thoughts1 }\OtherTok{\textless{}{-}} \FunctionTok{findThoughts}\NormalTok{(productsFit, }\AttributeTok{texts =}\NormalTok{ shortdoc, }\AttributeTok{n =} \DecValTok{2}\NormalTok{,}
\AttributeTok{topics =} \DecValTok{4}\NormalTok{)}\SpecialCharTok{$}\NormalTok{docs[[}\DecValTok{1}\NormalTok{]]}

\CommentTok{\# Find 2 documents associated with topic 5}
\NormalTok{thoughts2 }\OtherTok{\textless{}{-}} \FunctionTok{findThoughts}\NormalTok{(productsFit, }\AttributeTok{texts =}\NormalTok{ shortdoc, }\AttributeTok{n =} \DecValTok{2}\NormalTok{,}
\AttributeTok{topics =} \DecValTok{5}\NormalTok{)}\SpecialCharTok{$}\NormalTok{docs[[}\DecValTok{1}\NormalTok{]]}

\CommentTok{\# Plot the results}
\FunctionTok{par}\NormalTok{(}\AttributeTok{mfrow =} \FunctionTok{c}\NormalTok{(}\DecValTok{1}\NormalTok{, }\DecValTok{2}\NormalTok{), }\AttributeTok{mar =} \FunctionTok{c}\NormalTok{(}\FloatTok{0.5}\NormalTok{, }\FloatTok{0.5}\NormalTok{, }\DecValTok{1}\NormalTok{, }\FloatTok{0.5}\NormalTok{))}
\FunctionTok{plotQuote}\NormalTok{(thoughts1, }\AttributeTok{width =} \DecValTok{30}\NormalTok{, }\AttributeTok{main =} \StringTok{"Topic 4"}\NormalTok{)}
\FunctionTok{plotQuote}\NormalTok{(thoughts2, }\AttributeTok{width =} \DecValTok{30}\NormalTok{, }\AttributeTok{main =} \StringTok{"Topic 5"}\NormalTok{)}
\end{Highlighting}
\end{Shaded}

\begin{figure}[H]

\centering{

\includegraphics{7.Chapter_files/figure-pdf/fig-stm5-1.pdf}

}

\caption{\label{fig-stm5}Topic representative document excerpts}

\end{figure}%

The results in Figure~\ref{fig-stm5} show excerpts of the
first 200 characters of the most representative two documents for topics
4 and 5. These give us an insight into what the topics are about,
providing a qualitative view which complements the model's quantitative
approach.

\subsection{Estimating metadata-topic
relationships}\label{estimating-metadata-topic-relationships}

The main function in the \texttt{stm} package for evaluating the link
between topics and metadata is \texttt{estimateEffect()}.
For \(K\) topics, it identifies their relationships and the associated
uncertainty by employing, generally, the topical-prevalence model used
in estimating STM. The topical-prevalence model is similar to a
regression model, the topic proportions represent the dependent variable,
while metadata (e.g., \texttt{Brand}, \texttt{Review\_date} and
\texttt{Rating}) are the independent variables which explain how the
prevalence of topics varies across documents. As an input, the function can take selected topics and a link formula for the metadata. There
are several options for estimating this link function. To speed up
processing time, you can choose \texttt{uncertainty\ =\ "None"}, but
this will not give the estimation uncertainty, or, you can use \texttt{"Global"},
which does consider the estimation uncertainty of the topic proportions
(Roberts, Stewart, and Tingley 2019). The estimated relationships can be
useful in interpreting and validating the topics (Grimmer and
Stewart 2013). First, we need to transform the \texttt{Brand} column
into a categorical variable, called a factor in R. Then we need to
estimate the effect of the covariates, in this case we have the
\texttt{Brand}, the \texttt{Review\_date}, and the product's
\texttt{Rating} provided in a document ranging from 1 to 5 in the
\texttt{productsFit} topic model. The function \texttt{s()} applies a
spline to the \texttt{Review\_date} variable. This allows us to capture
non-linear trends over time. The role of a spline is to create a smooth
curve that models variations in topic prevalence as a continuous,
non-linear function applied to a variable, in our case
\texttt{Review\_date}. The \texttt{summary()} function gives a summary
of the estimated effect of the topic.

\begin{Shaded}
\begin{Highlighting}[]
\CommentTok{\# Convert "Brand" to a categorical variable}
\NormalTok{out}\SpecialCharTok{$}\NormalTok{meta}\SpecialCharTok{$}\NormalTok{Brand }\OtherTok{\textless{}{-}} \FunctionTok{as.factor}\NormalTok{(out}\SpecialCharTok{$}\NormalTok{meta}\SpecialCharTok{$}\NormalTok{Brand)}

\CommentTok{\# Estimate the effect of the covariates on topics}
\NormalTok{prep }\OtherTok{\textless{}{-}} \FunctionTok{estimateEffect}\NormalTok{(}\DecValTok{1}\SpecialCharTok{:}\DecValTok{5} \SpecialCharTok{\textasciitilde{}}\NormalTok{ Brand }\SpecialCharTok{+}\FunctionTok{s}\NormalTok{(Review\_date)}\SpecialCharTok{+}\NormalTok{Rating, productsFit, }\AttributeTok{meta =}\NormalTok{ out}\SpecialCharTok{$}\NormalTok{meta, }\AttributeTok{uncertainty =} \StringTok{"Global"}\NormalTok{)}

\CommentTok{\# Summarize the estimated effect on topic 1}
\FunctionTok{summary}\NormalTok{(prep, }\AttributeTok{topics =} \DecValTok{1}\NormalTok{)}
\end{Highlighting}
\end{Shaded}

\begin{verbatim}

Call:
estimateEffect(formula = 1:5 ~ Brand + s(Review_date) + Rating, 
    stmobj = productsFit, metadata = out$meta, uncertainty = "Global")


Topic 1:

Coefficients:
                   Estimate Std. Error t value Pr(>|t|)    
(Intercept)       0.5055728  0.0171859  29.418   <2e-16 ***
BrandApple        0.0004276  0.0088018   0.049    0.961    
s(Review_date)1   0.0004648  0.0295211   0.016    0.987    
s(Review_date)2  -0.0030919  0.0190860  -0.162    0.871    
s(Review_date)3   0.0011639  0.0231625   0.050    0.960    
s(Review_date)4  -0.0042774  0.0191283  -0.224    0.823    
s(Review_date)5   0.0040100  0.0240889   0.166    0.868    
s(Review_date)6  -0.0126216  0.0205136  -0.615    0.538    
s(Review_date)7   0.0036750  0.0221737   0.166    0.868    
s(Review_date)8   0.0123759  0.0261013   0.474    0.635    
s(Review_date)9  -0.0103679  0.0260787  -0.398    0.691    
s(Review_date)10  0.0284871  0.0238802   1.193    0.233    
Rating           -0.0762326  0.0017742 -42.967   <2e-16 ***
---
Signif. codes:  0 '***' 0.001 '**' 0.01 '*' 0.05 '.' 0.1 ' ' 1
\end{verbatim}

Analyzing the output, we can come to the following conclusions: The
intercept has a significant high t-value, this means that the base level
of the topic proportion is markedly different from zero. The coefficient
for brand Apple is very small (0.0004276) and has a high p-value (p =
0.961, so is statistically insignificant), meaning that Apple is
not strongly associated with Topic 1 when compared with Amazon. The
coefficients for the spline function applied to \texttt{Review\_date}
are all statistically insignificant, so the date of the review does not
influence Topic 1. The coefficient for \texttt{Rating} is negative
(-0.0770018) and highly significant (p\textless2e-16), meaning that
higher ratings are associated with a lower prevalence of Topic 1. In other
words, Topic 1 is more prevalent in reviews with lower ratings, this can
indicate the topic is related to user unhappiness. This information can be
used by the companies to  focus on specific themes and to address
the concerns of users.

In marketing, comparing brand associations is essential for
brand management. The change in topic proportions as a function of
associated brands can be evaluated using \texttt{method="difference"}.
In our case, this means taking the words from the score method value in
\texttt{productsFit} and plotting the change in topic proportion between the
two brands, Apple and Amazon. The \texttt{Score} method captures the
frequent words in a topic while penalizing their prevalence across all the
topics. This is ideal for identifying specific terms within each topic.

\begin{Shaded}
\begin{Highlighting}[]
\CommentTok{\# Using apply to concatenate labels for each row}
\NormalTok{custom\_labels }\OtherTok{\textless{}{-}} \FunctionTok{apply}\NormalTok{(lT}\SpecialCharTok{$}\NormalTok{score, }\DecValTok{1}\NormalTok{, }\ControlFlowTok{function}\NormalTok{(x) }\FunctionTok{paste}\NormalTok{(x, }\AttributeTok{collapse =} \StringTok{", "}\NormalTok{))}

\CommentTok{\# Plot the results}
\FunctionTok{plot}\NormalTok{(prep, }\AttributeTok{covariate =} \StringTok{"Brand"}\NormalTok{, }\AttributeTok{topics =} \FunctionTok{c}\NormalTok{(}\DecValTok{1}\NormalTok{, }\DecValTok{2}\NormalTok{, }\DecValTok{3}\NormalTok{, }\DecValTok{4}\NormalTok{, }\DecValTok{5}\NormalTok{),}
     \AttributeTok{model =}\NormalTok{ productsFit, }\AttributeTok{method =} \StringTok{"difference"}\NormalTok{, }\AttributeTok{cov.value1 =} \StringTok{"Apple"}\NormalTok{,}
     \AttributeTok{cov.value2 =} \StringTok{"Amazon"}\NormalTok{, }\AttributeTok{xlab =} \StringTok{"More Amazon ... More Apple"}\NormalTok{,}
     \AttributeTok{xlim =} \FunctionTok{c}\NormalTok{(}\SpecialCharTok{{-}}\FloatTok{0.5}\NormalTok{, }\FloatTok{0.5}\NormalTok{), }\AttributeTok{labeltype =} \StringTok{"custom"}\NormalTok{, }\AttributeTok{custom.labels =}\NormalTok{ custom\_labels)}
\end{Highlighting}
\end{Shaded}

\begin{figure}[H]

\centering{

\includegraphics{7.Chapter_files/figure-pdf/fig-stm7-1.pdf}

}

\caption{\label{fig-stm7}Effect of Amazon vs.~Apple}

\end{figure}%

From the plot in Figure~\ref{fig-stm7}, we can see that topics 1, 2, 3,
and 5, are more strongly associated with Amazon while Topic 4 is more
associated with Apple. It is evident that topic 1, the most negative, is defined by words related to customer service, at least for
some users. Marketers can make use of the information which is now
easy to see: The Kindle Fire might be an important Christmas gift. These
insights can us to help understand how specific brands dominate certain
themes.

There are variables which can be treated as continuous, such as the
\texttt{Review\_date}. Plotting the relationship
between the time (Year 2012) and its effect on topic 4 would yield
interesting results, see Figure~\ref{fig-stm8}. Remember, the variable
\texttt{prep} contains the estimated effects.

\begin{Shaded}
\begin{Highlighting}[]
\CommentTok{\# Plot the effect of \textquotesingle{}Review\_date\textquotesingle{} on topic 4 }
\FunctionTok{plot}\NormalTok{(prep, }\StringTok{"Review\_date"}\NormalTok{, }\AttributeTok{method =} \StringTok{"continuous"}\NormalTok{, }\AttributeTok{topics =} \DecValTok{4}\NormalTok{, }\AttributeTok{model =}\NormalTok{ productsFit, }\AttributeTok{printlegend =} \ConstantTok{FALSE}\NormalTok{, }\AttributeTok{xaxt =} \StringTok{"n"}\NormalTok{, }
\AttributeTok{xlab=}\StringTok{"Time (2012)"}\NormalTok{)}
\CommentTok{\# Create a sequence of dates incremented by month.}
\NormalTok{monthseq }\OtherTok{\textless{}{-}} \FunctionTok{seq}\NormalTok{(}\AttributeTok{from =} \FunctionTok{as.Date}\NormalTok{(}\StringTok{"2012{-}01{-}01"}\NormalTok{), }\AttributeTok{to =} \FunctionTok{as.Date}\NormalTok{(}\StringTok{"2012{-}07{-}11"}\NormalTok{), }\AttributeTok{by =} \StringTok{"month"}\NormalTok{)}

\CommentTok{\# Extract the month name}
\NormalTok{monthnames }\OtherTok{\textless{}{-}} \FunctionTok{months}\NormalTok{(monthseq)}

\CommentTok{\# Add custom x{-}axis labels to the plot}
\FunctionTok{axis}\NormalTok{(}\DecValTok{1}\NormalTok{,}\AttributeTok{at =} \FunctionTok{as.numeric}\NormalTok{(monthseq) }\SpecialCharTok{{-}} \FunctionTok{min}\NormalTok{(}\FunctionTok{as.numeric}\NormalTok{(monthseq)), }\AttributeTok{labels =}\NormalTok{ monthnames)}
\end{Highlighting}
\end{Shaded}

\begin{figure}[H]

\centering{

\includegraphics{7.Chapter_files/figure-pdf/fig-stm8-1.pdf}

}

\caption{\label{fig-stm8}The Effect of `Review\_date' on Topic 4}

\end{figure}%

Because there is a bump starting from March and April and going
down in May, we might imagine Apple held a special event in 2012. And indeed, on March 7 2012,
Apple hosted an event, announcing it with the tagline
 `We have something you really have to see. And touch'', where they presented the
3rd generation iPad, a model with a Retina display and a
dual-core Apple A5X processor with quad-core graphics. So, the timeline of this topic
shows that something exceptional for Apple
happened during that time interval.

\subsection{Topical content.}\label{topical-content.}

To check the influence of a topical content covariate, we need to
understand how the vocabulary used to describe a particular topic varies
as a function of this. It is achieved by fitting the STM
with a content variable such as \texttt{Brand}. When we include a content covariate, we examine how the
words associated with each topic differ depending on the value of the
covariate (e.g., Apple vs.~Amazon for \texttt{Brand}). We do
this by fitting a completely new model and thus the new topics may
differ from the previous example where no content covariate was
used.

\begin{Shaded}
\begin{Highlighting}[]
\CommentTok{\# Fit the model with the Brand content}
\NormalTok{dataContent }\OtherTok{\textless{}{-}} \FunctionTok{stm}\NormalTok{(out}\SpecialCharTok{$}\NormalTok{documents, out}\SpecialCharTok{$}\NormalTok{vocab, }\AttributeTok{K =} \DecValTok{5}\NormalTok{,}
 \AttributeTok{prevalence =}\SpecialCharTok{\textasciitilde{}}\NormalTok{ Brand }\SpecialCharTok{+} \FunctionTok{s}\NormalTok{(Review\_date), }\AttributeTok{content =}\SpecialCharTok{\textasciitilde{}}\NormalTok{ Brand,}
 \AttributeTok{max.em.its =} \DecValTok{75}\NormalTok{, }\AttributeTok{data =}\NormalTok{ out}\SpecialCharTok{$}\NormalTok{meta, }\AttributeTok{init.type =} \StringTok{"Spectral"}\NormalTok{,}
 \AttributeTok{verbose =} \ConstantTok{FALSE}\NormalTok{)}
\end{Highlighting}
\end{Shaded}

\subsection{Visualizing the influence of topical
content}\label{visualizing-the-influence-of-topical-content}

To visualize the influence of a topical content covariate (in our case
Brand), we can use the plot method for the newly created STM object with
the option = ``perspectives''. This visualization compares the
vocabulary used within a specific topic (e.g., Topic 4) with different
values of the covariate (e.g., Apple vs.~Amazon for \texttt{Brand}). This
function displays which words within a topic are more associated with
one covariate value than another. In the vocabulary differences
by brand for topic 4, that of Apple, the terms identified
by the \texttt{covariate} output show how the vocabulary
varies between Apple and Amazon.

\begin{Shaded}
\begin{Highlighting}[]
\CommentTok{\# Plot vocabulary differences by Brand for Topic 4}
\FunctionTok{plot}\NormalTok{(dataContent, }\AttributeTok{type =} \StringTok{"perspectives"}\NormalTok{, }\AttributeTok{topics =} \DecValTok{4}\NormalTok{)}
\end{Highlighting}
\end{Shaded}

\begin{figure}[H]

\centering{

\includegraphics{7.Chapter_files/figure-pdf/fig-stm10-1.pdf}

}

\caption{\label{fig-stm10}Brand-Specific Vocabulary in Topic 4}

\end{figure}%

The larger words on each side of the plot, see Figure~\ref{fig-stm10},
are those most strongly associated with the covariate level (Apple or
Amazon) for Topic 4.

\begin{Shaded}
\begin{Highlighting}[]
\CommentTok{\# Label the top words for Topic 4}
\NormalTok{(lT2 }\OtherTok{\textless{}{-}} \FunctionTok{labelTopics}\NormalTok{(dataContent, }\DecValTok{4}\NormalTok{))}
\end{Highlighting}
\end{Shaded}

\begin{verbatim}
Topic Words:
 Topic 4: thinner, processor, macbook, nexus, gen, higher, acer 
 
 Covariate Words:
 Group Amazon: small, owner, capabl, download, seri, say, amazon 
 Group Apple: ibook, imac, appl, wherea, industri, profession, frame 
 
 Topic-Covariate Interactions:
 Topic 4, Group Amazon: green, pad, creativ, clever, mirror, damn, window 
 Topic 4, Group Apple: seller, packag, condit, scratch, citi, human, durabl 
 
\end{verbatim}

From inspecting the visualization we see that for Apple the emphasis is on
words such as, ibook, imac, appl, wherea, industri, profession, and frame, while for Amazon it is on small, owner, capabl, download, seri,
say, and amazon. There are obvious distancing patterns for the vocabulary of Topic 4
based on the \texttt{Brand} covariate. For Amazon, words like
``small'', ``owner'', and ``capable'' are prominent, perhaps reflecting the tablet's
usability and and the experience of owners. For Apple, terms like
ibook, imac, or profession, underline the products' branding and professional
associations. Words such as ``creativity'' and
adaptability (through to  words like ``green'' and ``clever'') show the interactive nature of the Amazon product,
while the Apple reviews highlight the durability of the products and also
the condition of the products, related the sellers and packaging. Insights which might be of use here for Amazon could be the usability of
the product, ownership experience, and high degree of versatility, appealing to a
broad audience, while for Apple, branding and professional associations
appeal to consumers who are industry-focused.

\subsection{Comparing words across two
topics}\label{comparing-words-across-two-topics}

We can also use the \texttt{plot()} function to check the contrast in
words across two topics, see Figure~\ref{fig-stm11}.

\begin{Shaded}
\begin{Highlighting}[]
\CommentTok{\# Plot the contrast in words across Topics 4 and 5}
\FunctionTok{plot}\NormalTok{(dataContent, }\AttributeTok{type =} \StringTok{"perspectives"}\NormalTok{, }\AttributeTok{topics =} \FunctionTok{c}\NormalTok{(}\DecValTok{5}\NormalTok{, }\DecValTok{4}\NormalTok{))}
\end{Highlighting}
\end{Shaded}

\begin{figure}[H]

\centering{

\includegraphics{7.Chapter_files/figure-pdf/fig-stm11-1.pdf}

}

\caption{\label{fig-stm11}Contrast in words across Topics 4 and 5}

\end{figure}%

To further explore the content, we can create a word cloud for a
specific topic, see Figure~\ref{fig-stm12}.

\begin{Shaded}
\begin{Highlighting}[]
\CommentTok{\# Create a word cloud for Topic 4}
\FunctionTok{cloud}\NormalTok{(dataContent, }\AttributeTok{topic =} \DecValTok{4}\NormalTok{)}
\end{Highlighting}
\end{Shaded}

\begin{figure}[H]

\centering{

\includegraphics{7.Chapter_files/figure-pdf/fig-stm12-1.pdf}

}

\caption{\label{fig-stm12}Wordcloud, Topic 4}

\end{figure}%

Topic correlations can be investigated too. These can identify groups of
topics which tend to appear together in documents, possibly revealing
latent structures. The correlations can be visualized after calculating
the topic correlation matrix using the function \texttt{topicCorr()}.

\begin{Shaded}
\begin{Highlighting}[]
\CommentTok{\# Calculate and plot topic correlations}
\NormalTok{mod.out.corr }\OtherTok{\textless{}{-}} \FunctionTok{topicCorr}\NormalTok{(dataContent)}
\FunctionTok{plot}\NormalTok{(mod.out.corr)}
\end{Highlighting}
\end{Shaded}

\begin{figure}[H]

\centering{

\includegraphics{7.Chapter_files/figure-pdf/fig-stm13-1.pdf}

}

\caption{\label{fig-stm13}Topic Correlations}

\end{figure}%

To interpret the topic correlation plot we need to define its
components. Each node in the plot represents a topic, each edge
signifies the correlation between two topics. Positive correlations are
typically depicted with solid lines, while negative correlations are
often shown with dashed lines. Thicker lines indicate stronger
correlations. Groups of tightly connected nodes may form clusters,
showing topics which frequently occur together in documents. In our example,
we can see that topics 4 and 2, as well as topics 5 and 3, are negatively
correlated, whereas topic 1 appears to be more detached from the the other
topics.

\subsection{Exploring topic interactions and visualizations in
STM}\label{exploring-topic-interactions-and-visualizations-in-stm}

It is also possible to explore complex relationships between covariates and topics in
STM. To understand how the interactions between covariates can
reveal how one variable can moderate the effect of another, we can estimate a
model including a term which interacts with \texttt{Brand} and
\texttt{Review\_date}.

\begin{Shaded}
\begin{Highlighting}[]
\NormalTok{tabletInteraction }\OtherTok{\textless{}{-}} \FunctionTok{stm}\NormalTok{(out}\SpecialCharTok{$}\NormalTok{documents, out}\SpecialCharTok{$}\NormalTok{vocab, }\AttributeTok{K =} \DecValTok{5}\NormalTok{,}
\AttributeTok{prevalence =}\SpecialCharTok{\textasciitilde{}}\NormalTok{ Brand }\SpecialCharTok{*}\NormalTok{ Review\_date, }\AttributeTok{max.em.its =} \DecValTok{75}\NormalTok{, }\AttributeTok{data =}\NormalTok{ out}\SpecialCharTok{$}\NormalTok{meta,}
\AttributeTok{init.type =} \StringTok{"Spectral"}\NormalTok{, }\AttributeTok{verbose =} \ConstantTok{FALSE}\NormalTok{)}
\end{Highlighting}
\end{Shaded}

We have re-estimated the model to consider the interaction between the brand
and the day, and transformed this linearly. Using the function
\texttt{estimateEffect()} we have included the same interaction and prepared
these objects for plotting, showing the interaction over time.

\begin{Shaded}
\begin{Highlighting}[]
\NormalTok{prep }\OtherTok{\textless{}{-}} \FunctionTok{estimateEffect}\NormalTok{(}\FunctionTok{c}\NormalTok{(}\DecValTok{4}\NormalTok{) }\SpecialCharTok{\textasciitilde{}}\NormalTok{ Brand }\SpecialCharTok{*}\NormalTok{ Review\_date, tabletInteraction, }\AttributeTok{meta =}\NormalTok{ out}\SpecialCharTok{$}\NormalTok{meta, }\AttributeTok{uncertainty =} \StringTok{"None"}\NormalTok{)}

\CommentTok{\# Plot interaction effect for Apple}
\FunctionTok{plot}\NormalTok{(prep, }\AttributeTok{covariate =} \StringTok{"Review\_date"}\NormalTok{, }\AttributeTok{model =}\NormalTok{ tabletInteraction,}
  \AttributeTok{method =} \StringTok{"continuous"}\NormalTok{, }\AttributeTok{xlab =} \StringTok{"Days"}\NormalTok{, }\AttributeTok{moderator =} \StringTok{"Brand"}\NormalTok{,}
  \AttributeTok{moderator.value =} \StringTok{"Apple"}\NormalTok{, }\AttributeTok{linecol =} \StringTok{"blue"}\NormalTok{, }\AttributeTok{ylim =} \FunctionTok{c}\NormalTok{(}\DecValTok{0}\NormalTok{, }\FloatTok{0.60}\NormalTok{),}
  \AttributeTok{printlegend =} \ConstantTok{FALSE}\NormalTok{)}

\CommentTok{\# Add interaction effect for Amazon}
\FunctionTok{plot}\NormalTok{(prep, }\AttributeTok{covariate =} \StringTok{"Review\_date"}\NormalTok{, }\AttributeTok{model =}\NormalTok{ tabletInteraction,}
  \AttributeTok{method =} \StringTok{"continuous"}\NormalTok{, }\AttributeTok{xlab =} \StringTok{"Days"}\NormalTok{, }\AttributeTok{moderator =} \StringTok{"Brand"}\NormalTok{,}
  \AttributeTok{moderator.value =} \StringTok{"Amazon"}\NormalTok{, }\AttributeTok{linecol =} \StringTok{"red"}\NormalTok{, }\AttributeTok{add =} \ConstantTok{TRUE}\NormalTok{,}
  \AttributeTok{printlegend =} \ConstantTok{FALSE}\NormalTok{)}

\FunctionTok{legend}\NormalTok{(}\DecValTok{0}\NormalTok{, }\FloatTok{0.4}\NormalTok{, }\FunctionTok{c}\NormalTok{(}\StringTok{"Apple"}\NormalTok{, }\StringTok{"Amazon"}\NormalTok{), }\AttributeTok{lwd =} \DecValTok{2}\NormalTok{,}
  \AttributeTok{col =} \FunctionTok{c}\NormalTok{(}\StringTok{"blue"}\NormalTok{, }\StringTok{"red"}\NormalTok{))}
\end{Highlighting}
\end{Shaded}

\begin{figure}[H]

\centering{

\includegraphics{7.Chapter_files/figure-pdf/fig-stm15-1.pdf}

}

\caption{\label{fig-stm15}Brand Interaction over time for Topic 4}

\end{figure}%

The plot in Figure~\ref{fig-stm15} shows the prevalence of Topic 4 as a
linear function of time, with the brand set to either Apple or Amazon.
Other variables in the model are held at their sample medians. From the
graph, it is clear that Apple is discussed widely in Topic 4,
slightly increasing over time, while Amazon is underrepresented and
discussed without much variation over time.

Again, we can explore the topic correlations using the
\texttt{tabletInteraction} model. Positive correlations indicate that
both topics are likely to be discussed within the same document.

\begin{Shaded}
\begin{Highlighting}[]
\CommentTok{\# Calculate the topic correlations}
\NormalTok{mod.out.corr }\OtherTok{\textless{}{-}} \FunctionTok{topicCorr}\NormalTok{(tabletInteraction)}

\CommentTok{\# View the correlation matrix}
\NormalTok{mod.out.corr}\SpecialCharTok{$}\NormalTok{cor}
\end{Highlighting}
\end{Shaded}

\begin{verbatim}
           [,1]        [,2]        [,3]        [,4]        [,5]
[1,]  1.0000000 -0.14830514 -0.45067563 -0.16832829 -0.32473195
[2,] -0.1483051  1.00000000 -0.46675733 -0.03698267 -0.30964228
[3,] -0.4506756 -0.46675733  1.00000000 -0.22282036  0.01811152
[4,] -0.1683283 -0.03698267 -0.22282036  1.00000000 -0.35590402
[5,] -0.3247319 -0.30964228  0.01811152 -0.35590402  1.00000000
\end{verbatim}

\begin{Shaded}
\begin{Highlighting}[]
\CommentTok{\# Plot the topic correlation network}
\FunctionTok{plot}\NormalTok{(mod.out.corr)}
\end{Highlighting}
\end{Shaded}

\begin{figure}[H]

\centering{

\includegraphics{7.Chapter_files/figure-pdf/fig-stm16-1.pdf}

}

\caption{\label{fig-stm16}Topic correlation using the tabletInteraction
model}

\end{figure}%

Unlike in Figure~\ref{fig-stm13} that shows the influence of the brand
covariate, in Figure~\ref{fig-stm16} we can see the general
co-occurrence of the topics and their opposition across the dataset. The negative
correlation between Topics 3 and 5 highlights a divergence in their
thematic content, while the independence of Topics 1, 2, and 4 indicate
that these topics cover distinct, non-overlapping themes.

Given the significant number of visualizations and information that can
be extracted from the data, this paves the way for further exploration,
for example, in the exercises at the end of this chapter.

\subsection{Determining the number of topics using
STM.}\label{determining-the-number-of-topics-with-stm.}

As specified in the \texttt{stm} manual (Roberts,
Stewart, and Tingley 2019), there is no clear way to decide which is
the best number of topics for a given corpus (Grimmer and Stewart 2013).
While the \texttt{stm()} function accepts a fixed number for modeling
the topics, \texttt{searchK()} uses a data-driven approach to
select the number of topics. The function performs several
automated tests to help the analyst choose the number of topics. For
example, it is possible to estimate an STM model of 3 to 7 topics and compare
the results for each of the criteria. The default initialization is spectral initialization because of its stability. This function will
also calculate a range of quantities of interest, such as average
exclusivity, semantic coherence, held-out likelihood, residuals, bound,
lower bound, and EM iterations. Exclusivity generally increases with the
number of topics (K), indicating the presence of more distinct topics. Semantic
coherence decreases with (K), with a higher degree of coherence suggesting more
meaningful topics and some difficulty in interpreting them. The
held-out likelihood becomes less negative as (K) increases, with higher
values being better. Residuals decrease with increasing (K), with lower
residuals being preferable. The bound (or lower bound) increases with
(K), where higher values are advantageous and lead to a better fit between the model
and the data. EM iterations vary, and, typically, fewer iterations
indicate better convergence, but this depends on other factors as well.

\begin{Shaded}
\begin{Highlighting}[]
\NormalTok{storage }\OtherTok{\textless{}{-}} \FunctionTok{searchK}\NormalTok{(out}\SpecialCharTok{$}\NormalTok{documents, out}\SpecialCharTok{$}\NormalTok{vocab, }\AttributeTok{K =} \DecValTok{3}\SpecialCharTok{:}\DecValTok{7}\NormalTok{, }
                   \AttributeTok{prevalence =}\SpecialCharTok{\textasciitilde{}}\NormalTok{ Brand }\SpecialCharTok{+} \FunctionTok{s}\NormalTok{(Review\_date), }
                   \AttributeTok{data =}\NormalTok{ out}\SpecialCharTok{$}\NormalTok{meta, }\AttributeTok{verbose =} \ConstantTok{FALSE}\NormalTok{)}
\end{Highlighting}
\end{Shaded}

When selecting the best (K) from a range of topic models, we should
consider the factors given above (i.e., exclusivity, semantic coherence,
held-out likelihood, residuals, bound, and the number of EM iterations).
Let us analyze these metrics for the new results:

\begin{verbatim}
  K   exclus    semcoh   heldout residual    bound   lbound em.its
1 3 8.016233 -50.44074 -6.240349 2.916356 -3650009 -3650007     14
2 4 8.208482 -52.92883  -6.20944  2.83403 -3641386 -3641383     13
3 5 8.480603  -55.0643 -6.200306 2.782149 -3634169 -3634164     14
4 6 8.560167 -56.44334 -6.194746 2.729681 -3629368 -3629362     13
5 7 8.444968 -56.71841 -6.189823 2.692228 -3622360 -3622352     11
\end{verbatim}

The optimal number of topics, based on the metrics, appears to be 7. For example, it
has a high exclusivity compared to \(K=6\) , this means that
topics at \(K=7\) are more distinct than those at \(K=6\). Despite having the
lowest semantic coherence (the topics may be less susceptible to interpretation), it
performs well for the held-out likelihood, showing the least-negative
value (meaning that it generalizes unseen data better), residuals (with the best
fit to the data), and bound (being the best for both high and low). The low number of EM iterations indicates that the model is stable and not
overfitted. So \(K=7\) appears to offer a good balance between the
metrics if data fit and generalization are important for predictive
modeling or explanatory purposes, otherwise, it might be better to choose \(K=5\) for
more meaningful and easily understood topics if you prefer easier interpretation and
semantic coherence.

There are other ways to select the number of topics, for example using
LDA tuning. This might be a useful tool to help optimize topic-modeling
performance by selecting an optimal number of topics. By tuning the LDA
parameters we can create models that represent the topics in the
corpus better, allowing the analyst to interpret the results and gain insights which can lead to action. The \texttt{ldatuning} package (Nikita and Nikita
2016) allows for the analysis of some metrics which either show minimization
(Cao et al. 2009), or maximization (Griffiths and Steyvers 2004). We have
loaded the dataset to R, pre-processed it, then applied several
metrics, such as ``Griffiths2004'', ``CaoJuan2009'', ``Arun2010'', and a
method, such as ``Gibbs''.

\begin{Shaded}
\begin{Highlighting}[]
\FunctionTok{options}\NormalTok{(}\AttributeTok{seededlda\_threads =} \DecValTok{4}\NormalTok{)}

\CommentTok{\# Read product reviews from a CSV file into a dataframe}
\NormalTok{prodreviews }\OtherTok{\textless{}{-}} \FunctionTok{read.csv}\NormalTok{(}\StringTok{"data/Product Reviews.csv"}\NormalTok{)}


\CommentTok{\# Select the desired columns}
\NormalTok{data }\OtherTok{\textless{}{-}}\NormalTok{ prodreviews[, }\FunctionTok{c}\NormalTok{(}\StringTok{"Content"}\NormalTok{, }\StringTok{"Rating"}\NormalTok{, }\StringTok{"Brand"}\NormalTok{, }\StringTok{"Review\_date"}\NormalTok{)]}

\CommentTok{\# Transform the character content of Review\_date into date type}
\NormalTok{data}\SpecialCharTok{$}\NormalTok{Review\_date }\OtherTok{\textless{}{-}} \FunctionTok{mdy}\NormalTok{(data}\SpecialCharTok{$}\NormalTok{Review\_date)}

\CommentTok{\# Select the year 2011}
\NormalTok{data }\OtherTok{\textless{}{-}}\NormalTok{ data }\SpecialCharTok{\%\textgreater{}\%} \FunctionTok{filter}\NormalTok{(}\FunctionTok{year}\NormalTok{(Review\_date) }\SpecialCharTok{==} \DecValTok{2012}\NormalTok{)}

\CommentTok{\# Use the strftime to transform the date into a number (day of the year)}
\NormalTok{data}\SpecialCharTok{$}\NormalTok{Review\_date }\OtherTok{\textless{}{-}} \FunctionTok{as.integer}\NormalTok{(}\FunctionTok{strftime}\NormalTok{(data}\SpecialCharTok{$}\NormalTok{Review\_date, }\AttributeTok{format =}\StringTok{"\%j"}\NormalTok{))}

\CommentTok{\# Select the relevant rows (Amazon and Apple)}
\NormalTok{data }\OtherTok{\textless{}{-}} \FunctionTok{subset}\NormalTok{(data, Brand }\SpecialCharTok{==} \StringTok{"Amazon"} \SpecialCharTok{|}\NormalTok{ Brand }\SpecialCharTok{==} \StringTok{"Apple"}\NormalTok{)}

\CommentTok{\# do some cleaning here....}

\CommentTok{\# Create a text corpus from the \textquotesingle{}Content\textquotesingle{} column of the product reviews dataframe}
\NormalTok{corp }\OtherTok{\textless{}{-}} \FunctionTok{corpus}\NormalTok{(data, }\AttributeTok{text\_field =} \StringTok{"Content"}\NormalTok{)}
\NormalTok{dfmt }\OtherTok{\textless{}{-}} \FunctionTok{TermDocumentMatrix}\NormalTok{(corp)}

\FunctionTok{library}\NormalTok{(}\StringTok{"ldatuning"}\NormalTok{)}
\NormalTok{result }\OtherTok{\textless{}{-}} \FunctionTok{FindTopicsNumber}\NormalTok{(}
  \AttributeTok{dtm=}\NormalTok{dfmt,}
  \AttributeTok{topics =} \FunctionTok{seq}\NormalTok{(}\AttributeTok{from =} \DecValTok{3}\NormalTok{, }\AttributeTok{to =} \DecValTok{7}\NormalTok{, }\AttributeTok{by =} \DecValTok{1}\NormalTok{),}
  \AttributeTok{metrics =} \FunctionTok{c}\NormalTok{(}\StringTok{"Griffiths2004"}\NormalTok{, }\StringTok{"CaoJuan2009"}\NormalTok{, }\StringTok{"Arun2010"}\NormalTok{),}
  \AttributeTok{method =} \StringTok{"Gibbs"}\NormalTok{,}
  \AttributeTok{control =} \FunctionTok{list}\NormalTok{(}\AttributeTok{seed =} \DecValTok{42}\NormalTok{),}
  \AttributeTok{mc.cores =} \DecValTok{8}\NormalTok{L,}
  \AttributeTok{verbose =} \ConstantTok{TRUE}
\NormalTok{)}
\end{Highlighting}
\end{Shaded}

\begin{verbatim}
fit models... done.
calculate metrics:
  Griffiths2004... done.
  CaoJuan2009... done.
  Arun2010... done.
\end{verbatim}

\begin{Shaded}
\begin{Highlighting}[]
\FunctionTok{suppressWarnings}\NormalTok{(}\FunctionTok{FindTopicsNumber\_plot}\NormalTok{(result))}
\end{Highlighting}
\end{Shaded}

\begin{figure}[H]

\centering{

\includegraphics{7.Chapter_files/figure-pdf/fig-selectlda-1.pdf}

}

\caption{\label{fig-selectlda}Find Optimal Number of Topics}

\end{figure}%

From the plots (see Figure~\ref{fig-selectlda}), it is easy to
see that K=7 is the best number of topics, found by either
minimizing or maximizing the values, with both converging and producing similar results.

To acquire topics of high quality we need to consider several factors.
One of them is the quality and quantity of the corpus and its
pre-processing. Another is the variety of topics in the
corpus. The number of grams (n-grams), may also influence the
context and coherence of topics: Here we have worked with 1-grams but 2
or more grams would add value to the final results. Last but not least, the
fine-tuning of parameters will optimize the model's performance to achieve better
accuracy and greater insights.

\emph{Exercise starts}

\textbf{Exercise 1}: Other packages used for LDA topic modeling. Use
another package, such as \texttt{tm} from CRAN to perform LDA.

\textbf{Exercise 2}: A 2-gram approach for better insights. Use a 2-gram
(also known as a bi-gram) approach using the \texttt{seededlda} package and
check the results.

\textbf{Exercise 3}: Topic modeling with residuals. Use the topic
modeling with residuals and seeded words at sentence level using the
\texttt{seededlda} package. You can get more information on this package
by reading the package's manual.

\textbf{Exercise 4}. Bi-grams in the \texttt{stm}package. Use the
\texttt{stm} package with bi-grams for simple topic extraction.

\textbf{Exercise 5}. Covariates and bi-grams. Use the \texttt{stm}
package with different covariates and bi-grams to compare two brands.

\emph{Exercise ends}

\emph{Conclusion starts}

The chapter began with an introduction to topic modeling, showing the
simple mathematical formulas which power the latent Dirichlet allocation
(LDA) algorithm. It provided a simple example of how topics are derived and how
words in a corpus are assigned to specific topics. Next, we looked at a practical application of LDA using the same data set
and using the \texttt{seededlda} package. To give a more complete overview of the LDA, we looked at important examples, such as seeded LDA, which initializes topics around
keywords, or sentence-based LDA, which treats each sentence as a
distinct topic. The
chapter then explored another topic modeling technique, structural topic
model (STM), to expand the scope of topic modeling from a
marketing perspective. \emph{Conclusion ends}

